\documentclass[11pt,letterpaper]{article}
\usepackage[T1]{fontenc}
\usepackage[utf8]{inputenc}
\usepackage{amsmath,amsthm,mathtools}
\usepackage{newtxtext,newtxmath}
\usepackage[margin=1in,headheight=15pt]{geometry}
\usepackage{microtype}
\usepackage{enumitem,booktabs,array,longtable}
\usepackage{xcolor}
\usepackage{tikz-cd}
\usepackage[most]{tcolorbox}
\usepackage{fancyhdr}
\usepackage{hyperref}
\usepackage[nameinlink,noabbrev]{cleveref}
\definecolor{ink}{HTML}{293830}
\definecolor{soft}{HTML}{F1F3EC}
\hypersetup{colorlinks=true,linkcolor=ink,citecolor=ink,urlcolor=ink,
  pdftitle={Hahn--Tate Uniformization at Arbitrary Valuation Rank},
  pdfauthor={Research manuscript prepared with ChatGPT},
  pdfsubject={Surcomplex elliptic curves, Hahn summability, and torsion ramification}}
\pagestyle{fancy}
\fancyhf{}\fancyhead[L]{\small\scshape Hahn--Tate uniformization}
\fancyhead[R]{\small 22 September 2026}\fancyfoot[C]{\thepage}
\renewcommand{\headrulewidth}{0.3pt}
\setlength{\parskip}{3.5pt}
\setlength{\parindent}{1.1em}
\setlist{nosep,leftmargin=1.6em}
\setcounter{tocdepth}{2}
\newtheorem{theorem}{Theorem}[section]
\newtheorem{lemma}[theorem]{Lemma}
\newtheorem{proposition}[theorem]{Proposition}
\newtheorem{corollary}[theorem]{Corollary}
\theoremstyle{definition}
\newtheorem{definition}[theorem]{Definition}
\newtheorem{example}[theorem]{Example}
\newtheorem{remark}[theorem]{Remark}
\newcommand{\C}{\mathbb C}
\newcommand{\R}{\mathbb R}
\newcommand{\Q}{\mathbb Q}
\newcommand{\Z}{\mathbb Z}
\newcommand{\N}{\mathbb N}
\newcommand{\No}{\mathbf{No}}
\newcommand{\supp}{\operatorname{supp}}
\newcommand{\lc}{\operatorname{lc}}
\newcommand{\res}{\operatorname{res}}
\newcommand{\trop}{\operatorname{trop}}
\newcommand{\Gal}{\operatorname{Gal}}
\newcommand{\ord}{\operatorname{ord}}
\newcommand{\Span}{\operatorname{span}}
\newcommand{\m}{\mathfrak m}
\newcommand{\OO}{\mathcal O}
\newcommand{\Hq}{H_\alpha}
\newcommand{\Uq}{U_q}
\newcommand{\K}[1]{\C((t^{#1}))}
\newcommand{\sha}{\texttt{a3124af79f66b8b9c196d76b4cbc5ac3938907c4}}
\newcommand{\formal}{\textup{formal}}
\newcommand{\ds}{\displaystyle}
\newtcolorbox{statusbox}{colback=soft,colframe=ink,boxrule=.45pt,
  arc=1pt,left=9pt,right=9pt,top=7pt,bottom=7pt}

\begin{document}
\begin{titlepage}
\vspace*{0.4in}
{\large\scshape Surreal and surcomplex research\par}
\vspace{0.28in}
{\fontsize{29}{33}\selectfont\bfseries Hahn--Tate Uniformization\par
at Arbitrary Valuation Rank\par}
\vspace{0.22in}
{\Large Exact summability domains, surcomplex elliptic curves,\par
and torsion ramification\par}
\vspace{0.32in}
{\large Research manuscript prepared for Vladimir Reshetnikov\par}
\vspace{0.07in}
{22 September 2026\par}
\vspace{0.3in}
\begin{abstract}
\noindent
Let $K=\C((t^\Gamma))$ be a Hahn field, with no restriction on the rank
or divisibility of the set-sized ordered abelian group $\Gamma$. For
$q\in K$ of positive valuation $\alpha$, we identify the exact domain on
which the classical bilateral theta and Tate coordinate families are
strongly Hahn-summable: it is the multiplicative subgroup $U_q$ whose
valuations belong to the principal convex subgroup generated by $\alpha$.
Outside this domain, leading exponents themselves contain a strictly
decreasing sequence. Nevertheless, the resulting map gives an isomorphism
$U_q/q^\Z\simeq E_q(K)$ onto \emph{all} rational points of the Tate elliptic
curve. The proof combines a rank-one coarsening inside the principal convex
subgroup with unrestricted formal evaluation at infinitesimals for a second,
coarser valuation. It therefore does not mistake Hahn summation for
convergence at every surreal scale.

We also prove compatibility with exponent embeddings, an exact reduction
sequence, a valuation-circle quotient, support-preserving recovery of $q$
from $j$, and the formula
$K(E_q[n])=K_{\Gamma+\Z\alpha/n}$ for the full torsion field. Every elliptic
curve over $\No[i]$ with negative valuation of its $j$-invariant is covered,
using set-sized normal-form workspaces. Explicit rank-two and genuinely
surcomplex examples exhibit both the forbidden parameter scales and the
arbitrarily finer infinitesimals that the parametrization still captures.
\end{abstract}
\vspace{0.12in}
\begin{statusbox}
\textbf{Status and scope.}
This manuscript takes the \emph{new-theorems} alternative, not a claimed
solution of a named published open problem. The candidate original
contribution is the exact-domain, arbitrary-rank uniformization package and
its support-theoretic proof. Tate's rank-one theorem, theta identities, and
the classical torsion construction are credited as inputs. A targeted
literature and repository audit found no matching formulation, but does not
certify priority. The proofs have not been independently refereed or
formalized in Lean. Included exact computations check finite identities,
not the infinite arguments.
\end{statusbox}
\vfill
{\small\textbf{Keywords.} Surcomplex numbers; Hahn series; Tate curve;
convex valuation subgroup; strong summability; formal groups; ramification.}
\end{titlepage}

\tableofcontents
\newpage

\section{The problem and the principal results}\label{sec:intro}

The classical multiplicative uniformization of an elliptic curve uses a
parameter $u\in k^\times$ modulo powers of a small element $q$. Over a
complete real-valued non-Archimedean field, Tate's theorem supplies the
isomorphism $k^\times/q^\Z\simeq E_q(k)$ \cite{Tate,Silverman}.
Replacing $k$ by a Hahn field of arbitrary rank is not a harmless change
of coefficients. A positive value $\alpha=v(q)$ need not dominate all
other values by its integer multiples. The same bilateral series can then
have non-well-ordered support.

There are two different questions. Which parameters admit the literal
Hahn formulas? And, after excluding the inadmissible parameters, do those
formulas still reach every elliptic-curve point? The second question does
not follow from a support calculation. It is the central issue addressed
here.

Throughout, $\Gamma$ is a set-sized ordered abelian group, written
additively, and
\[
 K=K_\Gamma:=\K{\Gamma}.
\]
The coefficient field $\C$ has the trivial valuation. Every integer,
degree, and index $n\in\Z$ in a series is an \emph{ordinary finite}
integer, not a surreal integer. We use $v(t^\gamma)=\gamma$ and
$v(0)=+\infty$.

\begin{definition}[The admissible group]\label{def:domain}
For $q\in K^\times$ with $\alpha=v(q)>0$, put
\begin{align}
 H_\alpha&=\{\gamma\in\Gamma: -N\alpha\le\gamma\le N\alpha
                    \text{ for some }N\in\N\},\label{eq:H}\\
 U_q&=v^{-1}(H_\alpha)\subseteq K^\times.\label{eq:U}
\end{align}
Thus $H_\alpha$ is the smallest convex subgroup of $\Gamma$ containing
$\alpha$. We call $U_q$ the \emph{Hahn-admissible multiplicative domain}.
\end{definition}

Define the standard Tate coefficients
\begin{align}
 s_k(q)&=\sum_{n\ge1}\frac{n^kq^n}{1-q^n},\label{eq:sk}\\
 a_4(q)&=-5s_3(q),&
 a_6(q)&=-\frac{5s_3(q)+7s_5(q)}{12},\label{eq:ak}\\
 E_q&:\quad y^2+xy=x^3+a_4(q)x+a_6(q).\label{eq:curve}
\end{align}
The point at infinity is denoted $O$. These definitions are legitimate
strong Hahn sums; the discriminant is nonzero by \cref{prop:discriminant}.

\begin{theorem}[Exact-domain Hahn--Tate uniformization]\label{thm:main}
Let $q\in K_\Gamma$ satisfy $v(q)=\alpha>0$.
The bilateral family
\begin{equation}
 \Theta_q(u)=\sum_{n\in\Z}(-1)^nq^{n(n-1)/2}u^n\label{eq:theta}
\end{equation}
is strongly Hahn-summable if and only if $u\in U_q$.
For $u\notin q^\Z$, each of the two bilateral families defining
\begin{align}
 X_q(u)&=\sum_{n\in\Z}\frac{q^nu}{(1-q^nu)^2}-2s_1(q),\label{eq:X}\\
 Y_q(u)&=\sum_{n\in\Z}\frac{(q^nu)^2}{(1-q^nu)^3}+s_1(q)\label{eq:Y}
\end{align}
is strongly Hahn-summable if and only if $u\in U_q$.

The map
\begin{equation}
 \Phi_q(u)=
 \begin{cases}(X_q(u),Y_q(u)),&u\notin q^\Z,\\ O,&u\in q^\Z,
 \end{cases}\label{eq:phi}
\end{equation}
is a surjective group homomorphism $U_q\to E_q(K)$ with kernel $q^\Z$.
In particular,
\begin{equation}\boxed{\quad U_q/q^\Z\ \simeq\ E_q(K_\Gamma).\quad}\label{eq:main}
\end{equation}
No divisibility, countability, finite-rank, or real-valuedness assumption
on $\Gamma$ is needed.
\end{theorem}

The summability assertions are proved in \cref{sec:theta,sec:coordinates};
the group and surjectivity assertions are proved in \cref{sec:uniformization}.
The qualification ``the bilateral families'' matters: this is an exact
maximality theorem for specified expressions, not a claim that no different
function can be assigned outside $U_q$.

\begin{theorem}[Exact arithmetic and scalar extension]\label{thm:arithmetic-main}
For an integer $n\ge2$, let
$\Lambda=\Gamma+\Z\alpha/n\subseteq\Gamma\otimes_\Z\Q$.
Inside $K_{\Gamma\otimes\Q}$ one has
\begin{equation}
 K_\Gamma(E_q[n])=K_\Gamma(q^{1/n})=K_\Lambda,\qquad
 [K_\Lambda:K_\Gamma]=\ord(\alpha+n\Gamma)\le n.\label{eq:torsion-main}
\end{equation}
A single point $\Phi_q(q^{1/n})$ generates this field through its
coordinates. The extension is cyclic and totally ramified.

An ordered embedding of exponent groups preserves the Hahn formulas and
$\Phi_q$. In a larger Hahn field the exact parameter domain is again the
inverse image of the principal convex subgroup generated by the image of
$\alpha$; it need not be the whole multiplicative group.
\end{theorem}

The torsion description is classical in its Kummer origin. The point of
\cref{thm:arithmetic-main} is the explicit arbitrary-group ramification
formula and its compatibility with the domain theorem, not a claim to
have discovered Tate torsion.

\begin{corollary}[Surcomplex elliptic curves]\label{cor:surcomplex-main}
Fix the Conway normal-form valuation on $\No[i]$, with
$v(\omega^{-\gamma})=\gamma$. For every nonzero $q$ with $\alpha=v(q)>0$,
\begin{equation}
 \{u\in\No[i]^\times: |v(u)|\le N\alpha\text{ for some }N\in\N\}/q^\Z
       \ \simeq\ E_q(\No[i]).\label{eq:No-main}
\end{equation}
Every elliptic curve $E$ over $\No[i]$ with $v(j(E))<0$ is isomorphic to
an $E_q$ of this kind. The parameter $q$ is uniquely determined by $j(E)$.
These are class statements whose every instance and proof is realized in
a set-sized Hahn workspace.
\end{corollary}

\subsection{Why the result is not an entire-function substitution}

If $H_\alpha=\Gamma$, the parameter domain is all of $K^\times$.
For a nonzero real-valued group this always happens. If
$\Gamma=\R\oplus_{\mathrm{lex}}\R$ and $\alpha=(0,1)$, however,
$H_\alpha=\{0\}\oplus\R$. Elements of valuation $(1,0)$ or $(-1,0)$
are excluded. Nevertheless, \emph{all} $K$-points of $E_q$ are represented.
In particular the restriction does not mean that only points whose
coordinates have $q$-bounded valuations are reached.

The accepted parameter $u=1+h$, where $v(h)$ dominates every $N\alpha$,
already produces points with
\[
 X_q(1+h)\sim h^{-2},\qquad Y_q(1+h)\sim-h^{-3}.
\]
Their coordinates can be vastly larger than every integral power of
$q^{-1}$. Smallness of $u$ and smallness of $u-1$ are different questions.

\subsection{Relation to the repository and to the literature}\label{sec:position}

The repository audit was targeted: its document map, the root description,
the spectral report, and the rank-one analytic report were inspected, and
code search was used for elliptic/Tate terminology. The pinned revision
is \sha\ \cite{Repo}. The closest existing report works over the fixed
field $\C((t^\R))$ and develops Tate \emph{algebras}, disks, annuli, and a
partial-theta example \cite{RepoRankOne}. A Tate algebra is not a Tate
elliptic curve. The present paper neither reclaims that analytic
infrastructure nor repeats the partial-theta zero calculation.

Tate's original proof explicitly starts with a nontrivial real-valued
absolute value and proves uniformization over any complete field with
that absolute value \cite{Tate}. Modern treatments include
\cite{Silverman,BGR}; the 2026 work of Eterovi\'c--Vermeulen concerns
$p$-adic definability and Tate uniformization \cite{EV}. Higher-rank
analytic geometry already exists, notably the Hahn analytification of
Foster--Ranganathan \cite{FR}. Thus neither ``higher-rank geometry'' nor
``non-Archimedean elliptic functions'' is new here.

The potentially new package is more specific: an \emph{if and only if}
Hahn-support domain; a proof that this restricted domain nevertheless
uniformizes every rational point; and exact compatibility with coarse
reduction, scalar extension, and exponent-group arithmetic. No assertion
of bibliographic uniqueness follows from the targeted search. Related
formulations in formal, adic, or model-theoretic geometry may subsume
parts of the result.

\section{Hahn summation and the relevant order structure}\label{sec:hahn}

\subsection{Strong summability}

A Hahn series is a formal expression
\[
 a=\sum_{\gamma\in\Gamma}a_\gamma t^\gamma,
 \qquad a_\gamma\in\C,
\]
whose support $\supp(a)=\{\gamma:a_\gamma\ne0\}$ is well ordered in the
increasing order. For $a\ne0$, its valuation is the least support element.
A family $(a_i)_{i\in I}$ is \emph{strongly summable} when
$\bigcup_i\supp(a_i)$ is well ordered and each exponent occurs in the
support of only finitely many $a_i$. Its sum is defined coefficientwise.
Neither an order limit nor a limit in the full valuation topology is part
of this definition.

\begin{lemma}[Positive-support calculus]\label{lem:neumann}
Let $S\subseteq\Gamma_{>0}$ be well ordered. Its additive monoid $S^*$,
including $0$, is well ordered. Every element of $\Gamma$ has only
finitely many representations as a finite unordered sum of members of
$S$, counting multiplicities. Consequently, for finitely many Hahn
series $z_1,\dots,z_d$ of positive valuation, any formal power series
$F\in\C[[Z_1,\dots,Z_d]]$ has a well-defined strong evaluation $F(z)$.
Evaluation preserves sums, products, and composition when the inner
formal series have zero constant coefficient.
\end{lemma}

\begin{proof}
We recall the order-theoretic argument behind Neumann's lemma
\cite{Neumann}. The alphabet $S$, with its order, is well quasi-ordered.
Higman's lemma says that finite words over this alphabet are well
quasi-ordered by order-preserving subsequence embedding with each letter
allowed to increase \cite{Higman}. If a word embeds into another, the
sum of its letters does not exceed the sum of the other word. The
inequality is strict if a letter increases or a positive letter is added.
A descending sequence of sums would therefore contradict Higman's lemma.
An infinite collection of distinct nondecreasing words of the same sum
would also contradict it: an embedding between two such words with equal
sum forces the words to coincide. Sorting a word removes the finitely
many permutations of each finite multiset. This proves both assertions.

The union of the finitely many positive supports $\supp(z_i)$ is well
ordered. Expanding all monomials in $F(z)$ produces exponents in $S^*$.
The finiteness assertion bounds the word lengths and the decompositions
at a fixed exponent, and the number of variable labels is finite.
Hence every coefficient is a finite sum. The same argument for the
jointly indexed families proves multiplication and the stated composition
rule. This is the usual Hahn--Neumann calculus, not a new summation
principle.
\end{proof}

In particular, $(1-z)^{-1}=\sum_{n\ge0}z^n$ and
$(1+z)^r=\sum_{n\ge0}\binom rn z^n$ for $r\in\Q$ and $v(z)>0$.
These identities hold even when $n v(z)$ is bounded above in $\Gamma$.
A series with leading term $ct^\gamma$ is inverted by factoring it as
$ct^\gamma(1+\varepsilon)$ and applying this rule to $\varepsilon$.

\subsection{Normalization and two convex subgroups}

\begin{lemma}[Finite normalization]\label{lem:normalize}
For $\beta\in\Gamma$, the following are equivalent:
\begin{enumerate}[label=(\roman*)]
\item $\beta\in H_\alpha$;
\item there is an integer $m$ with $0\le\beta+m\alpha<\alpha$.
\end{enumerate}
Such an $m$ is unique. If $\beta\notin H_\alpha$, then either
$\beta>N\alpha$ for every $N\in\N$, or $\beta<-N\alpha$ for every
$N\in\N$.
\end{lemma}

\begin{proof}
If $|\beta|\le N\alpha$, the finite list of adjacent half-open intervals
$[k\alpha,(k+1)\alpha)$ for sufficiently many integers $k$ covers
$\beta$. Their interiors and endpoint convention give existence and
uniqueness. The converse is immediate. Convexity of $H_\alpha$ gives the
last alternative, or it follows directly by separating the two signs.
\end{proof}

There is another convex subgroup inside $H_\alpha$:
\begin{equation}
 H_\alpha^-=
 \{\gamma\in H_\alpha:n|\gamma|<\alpha\text{ for all }n\ge1\}.
 \label{eq:Hminus}
\end{equation}
It consists of the values infinitesimal relative to $\alpha$.
The quotient $H_\alpha/H_\alpha^-$ is nonzero and Archimedean.
Indeed, for every positive $\gamma\notin H_\alpha^-$ some multiple of
$\gamma$ is at least $\alpha$, and every element of $H_\alpha$ is bounded
by a multiple of $\alpha$. This proves the Archimedean comparison in the
quotient. H\"older's theorem embeds this quotient in $(\R,+)$; normalize
the embedding so that $\alpha$ has image $1$.

The two quotients serve different purposes. The quotient by
$H_\alpha^-$ produces a rank-one valuation \emph{inside} $K_{H_\alpha}$.
The quotient by $H_\alpha$ produces a coarser valuation on the \emph{whole}
field $K_\Gamma$, whose residue field is $K_{H_\alpha}$.
Interchanging them would destroy the proof.

\section{The exact domain of the bilateral theta series}\label{sec:theta}

\begin{theorem}[Theta domain]\label{thm:theta}
The family in \eqref{eq:theta} is strongly summable exactly for
$u\in U_q$. On that domain,
\begin{equation}
 \Theta_q(qu)=-u^{-1}\Theta_q(u).\label{eq:theta-functional}
\end{equation}
\end{theorem}

\begin{proof}
Write $\beta=v(u)$. The leading exponent of its $n$th summand is
\begin{equation}
 b_n=\frac{n(n-1)}2\alpha+n\beta.\label{eq:bn}
\end{equation}
Suppose first that $\beta\in H_\alpha$ and choose $N$ with
$|\beta|\le N\alpha$. The differences
$b_{n+1}-b_n=n\alpha+\beta$ show that the positive tail is eventually
strictly increasing. Writing $n=-k$ gives the same property for the tail
as $k\to+\infty$. Each tail eventually exceeds every prescribed element
of $H_\alpha$, by comparison with integer multiples of $\alpha$.
Together with the finitely many remaining indices, the set
$B=\{b_n:n\in\Z\}$ is well ordered. Each value of $b_n$ occurs at most
twice: if $b_n=b_m$ and $n\ne m$, then
$(n+m-1)\alpha+2\beta=0$.

Factor
$q=c_qt^\alpha(1+\varepsilon_q)$ and
$u=c_ut^\beta(1+\varepsilon_u)$ with positive-valuation errors.
Let $M$ be the monoid generated by their two supports. Every summand in
\eqref{eq:theta} has support in $b_n+M$, by \cref{lem:neumann}, including
negative powers of $u$. The sum of two well-ordered subsets of an ordered
abelian group is well ordered, with finitely many pairs giving any fixed
sum. Thus $B+M$ is well ordered. There are finitely many choices of $n$
at a fixed leading value and finitely many decompositions in $M$.
This proves strong summability of the original family.

If $\beta>N\alpha$ for all $N$, put $n=-k$. Then
\[
 b_{-(k+1)}-b_{-k}=(k+1)\alpha-\beta<0.
\]
These are actual nonzero leading coefficients, so their decreasing
exponents already obstruct well-ordering of the union of supports.
If $\beta<-N\alpha$ for all $N$, the sequence $b_n$ decreases for
$n\ge0$. This proves necessity without assuming that possible
cancellations between summands can repair an inadmissible family.

Finally, multiplication by $q$ shifts the index in a strongly summable
family. The coefficient equality
$q^n(-1)^nq^{n(n-1)/2}=-(-1)^{n+1}q^{n(n+1)/2}$ proves
\eqref{eq:theta-functional}.
\end{proof}

\begin{proposition}[Product and zeros]\label{prop:theta-product}
For $0\le v(u)<\alpha$, the Jacobi product is a Hahn identity:
\begin{equation}
 \Theta_q(u)=(1-u)\prod_{n\ge1}
          (1-q^n)(1-uq^n)(1-u^{-1}q^n).\label{eq:triple}
\end{equation}
The product means the strong sum over its finite subproducts after
expansion. The zeros on $U_q$ are exactly $q^\Z$, and all are simple
in the local external-variable sense.
\end{proposition}

\begin{proof}
All of $q$, $qu$, and $q/u$ have positive valuation. Their support monoid
licenses the expanded product and its regrouping. The classical Jacobi
triple-product identity, in this normalization, is an identity of the
corresponding formal Laurent families \cite{DLMF}; strong evaluation
therefore yields \eqref{eq:triple}. One can also compare coefficients of
$q$ after treating $u$ as a Laurent indeterminate and then use the same
support argument for specialization.

The product after the factor $1-u$ is $1$ plus a positive-valuation
series, hence a unit. Therefore the only zero in the normalized annulus
is $u=1$. The functional equation and \cref{lem:normalize} give all
zeros. At $u=1$ the derivative of the product is
\[
 \Theta_q'(1)=-\prod_{n\ge1}(1-q^n)^3\ne0.
\]
Reindexing transports simplicity to every $q^m$. This is differentiation
of a local power series in a variable; no derivation on the scalar field
is being imposed.
\end{proof}

\section{Tate coordinates, support certificates, and precision}\label{sec:coordinates}

\subsection{A form with positive tails}

Put
\[
 f(z)=\frac{z}{(1-z)^2},\qquad
 g(z)=\frac{z^2}{(1-z)^3},\qquad
 h(z)=\frac{z}{(1-z)^3}.
\]
The elementary identities $f(z^{-1})=f(z)$ and
$g(z^{-1})=-h(z)$ give, after normalizing $0\le\beta=v(u)<\alpha$,
\begin{align}
 X_q(u)&=f(u)+\sum_{n\ge1}\bigl(f(q^nu)+f(q^n/u)\bigr)-2s_1(q),
       \label{eq:Xpositive}\\
 Y_q(u)&=g(u)+\sum_{n\ge1}\bigl(g(q^nu)-h(q^n/u)\bigr)+s_1(q).
       \label{eq:Ypositive}
\end{align}
The minus sign and the unsquared numerator in the last tail of
\eqref{eq:Ypositive} are important. In contrast, the numerator in the
original bilateral formula \eqref{eq:Y} is squared.

At $v(u)=0$ we do \emph{not} expand the isolated expressions $f(u)$ and
$g(u)$ geometrically in $u$. They are finite rational operations, defined
unless $u=1$. Their denominators may be arbitrarily small.

\begin{theorem}[Coordinate domain and support certificate]\label{thm:coordinates}
For $u\notin q^\Z$, the families in \eqref{eq:X} and \eqref{eq:Y} are
strongly summable exactly when $u\in U_q$.
For $0\le v(u)<\alpha$ and $u\ne1$, let
\begin{equation}
 S=\supp(q)\cup\supp(qu)\cup\supp(q/u)\subseteq\Gamma_{>0}.
 \label{eq:S}
\end{equation}
Then
\begin{equation}
 \supp(X_q(u)-f(u)),\ \supp(Y_q(u)-g(u))
        \subseteq S^*\setminus\{0\}.\label{eq:certificate}
\end{equation}
The bilateral and positive-tail formulas agree.
\end{theorem}

\begin{proof}
The series $q^nu=q^{n-1}(qu)$ and $q^n/u=q^{n-1}(q/u)$ have positive
valuation for $n\ge1$. Expanding $f,g,h$ in those arguments and
expanding each term of $s_1$ gives exponents in $S^*\setminus\{0\}$.
The full family, including $n$ and every geometric-expansion index, is
strongly summable: at a fixed exponent, the Neumann finiteness theorem
bounds the total word length, hence also $n$ and the geometric index.
This proves \eqref{eq:certificate} and justifies every regrouping in
\eqref{eq:Xpositive}--\eqref{eq:Ypositive}. Multiplication of $u$ by a
power of $q$ merely reindexes the bilateral sums. Finite normalization
therefore proves sufficiency on $U_q$.

For necessity, suppose first $\beta>N\alpha$ for every $N$. All $q^nu$
have positive valuation. The leading exponents in the $X$ family are
$\beta+n\alpha$ and in the $Y$ family are $2\beta+2n\alpha$.
Both decrease strictly as $n\to-\infty$. If $\beta<-N\alpha$ for every
$N$, all $q^nu$ have negative valuation. The rational identities above
show that the leading exponents of both families are
$-\beta-n\alpha$, with a nonzero sign in the $Y$ case. These decrease
as $n\to+\infty$. In either case the original families fail the
well-ordering requirement.
\end{proof}

\begin{proposition}[Exact finite-tail bound]\label{prop:tail}
For normalized $u$, let $X^{[N]},Y^{[N]}$ retain the isolated rational
terms and the terms $1\le n\le N$ in the positive-tail formulas, and
replace $s_1$ by $\sum_{n=1}^N nq^n/(1-q^n)$. Then
\begin{equation}
 v(X_q-X^{[N]})\ge(N+1)\alpha-\beta,
 \qquad v(Y_q-Y^{[N]})\ge(N+1)\alpha-\beta.\label{eq:tail}
\end{equation}
Here $v(0)=+\infty$, so cancellations cause no exception.
\end{proposition}

\begin{proof}
Each omitted $f(q^n/u)$ or $h(q^n/u)$ has valuation $n\alpha-\beta$;
the other omitted coordinate terms have at least that valuation.
The omitted $s_1$ terms have valuation $n\alpha$. Strong summation and
the valuation inequality give the bounds.
\end{proof}

If a requested precision $\delta$ lies in $H_\alpha$, a finite $N$
certifies \eqref{eq:tail} beyond $\delta$. If $\delta>N\alpha$ for
all $N$, this certificate never reaches it. A valid exact Hahn sum has
therefore not been converted into a convergence theorem at all scales.

\subsection{The elliptic curve and its invariant differential}

\begin{proposition}[Discriminant and modular invariant]\label{prop:discriminant}
The coefficients in \eqref{eq:ak} are Hahn evaluations of integer
power series in $q$. The Weierstrass discriminant and $j$-invariant are
\begin{align}
 \Delta(q)&=q\prod_{n\ge1}(1-q^n)^{24},&v(\Delta(q))&=\alpha,
       \label{eq:Delta}\\
 j(q)&=q^{-1}+744+196884q+21493760q^2+\cdots,&v(j(q))&=-\alpha.
       \label{eq:j}
\end{align}
In particular $E_q$ is a nonsingular elliptic curve over $K_\Gamma$.
\end{proposition}

\begin{proof}
The formulas for $s_k$ are licensed by positive-support calculus.
The integer divisibility of $(5n^3+7n^5)/12$ follows by checking its
factors modulo $3$ and $4$, or from the standard Tate coefficient
identity. With $a_1=1$, $a_2=a_3=0$, the usual Weierstrass formulas give
\[
 b_2=1,\quad b_4=2a_4,\quad b_6=4a_6,\quad b_8=a_6-a_4^2,
 \qquad c_4=1-48a_4.
\]
The formal identities for $\Delta$ and $c_4^3/\Delta$ are classical
\cite{Tate,Silverman}. They are identities in integer formal power
series (Laurent series for $j$), so positive-support evaluation proves
them here. The product multiplying $q$ is a unit of leading coefficient
$1$. Hence $\Delta\ne0$, and the displayed valuations follow.
\end{proof}

The coordinate functions satisfy the curve equation and the group law;
a proof valid at every rank is given in \cref{sec:uniformization}. The
following local identity can already be checked term by term.

\begin{proposition}[Differential and inversion identities]\label{prop:differential}
On the local external-variable germs in $U_q$, away from poles,
\begin{equation}
 u\frac{dX_q}{du}=X_q+2Y_q,\qquad
 X_q(u^{-1})=X_q(u),\qquad
 Y_q(u^{-1})=-X_q(u)-Y_q(u).\label{eq:diff}
\end{equation}
Consequently the invariant differential pulls back as
\begin{equation}
 \Phi_q^*\left(\frac{dx}{2y+x}\right)=\frac{du}{u}.\label{eq:omega}
\end{equation}
The equation for differentials is interpreted in regular local charts
also where the displayed affine denominator vanishes.
\end{proposition}

\begin{proof}
The rational identity $z f'(z)=f(z)+2g(z)$ and the constants
$-2s_1+2s_1=0$ prove the first formula. Differentiating the positive
 tails preserves their support estimates locally, so the interchange is
licensed. The inversion formulas follow from $f(z^{-1})=f(z)$ and
$g(z^{-1})=-f(z)-g(z)$, followed by reindexing. Equation
\eqref{eq:omega} is the first identity rewritten as a differential.
It extends in regular charts since an invariant differential has no
pole on a smooth elliptic curve.
\end{proof}

Nothing in this proposition says that a chosen scalar derivation
$\partial:K\to K$ satisfies these formulas without the chain rule's
$q$-dependence. In particular, it makes no assertion about the
Berarducci--Mantova derivation of surreal scalars.

\section{Coarse infinitesimals admit unrestricted formal evaluation}\label{sec:coarse}

This section supplies the step that a bare appeal to ``completeness of
Hahn fields'' would miss. We shall evaluate formal series whose
coefficients have no uniform bound for the original valuation. A coarser
valuation makes this possible.

\subsection{The coarse residue field}

Let $H\subseteq\Gamma$ be any convex subgroup. The quotient
$\overline\Gamma=\Gamma/H$ is ordered. Define
\[
 w(a)=v(a)+H\quad(a\ne0),\qquad
 \OO_w=\{a:w(a)\ge0\},\qquad
 \m_w=\{a:w(a)>0\}.
\]
For $a\in\OO_w$, retaining exactly the terms with exponents in $H$
gives a residue homomorphism
\begin{equation}
 \res_H:\OO_w\longrightarrow K_H,\qquad
 \sum_\gamma a_\gamma t^\gamma\longmapsto
             \sum_{\gamma\in H}a_\gamma t^\gamma.\label{eq:resH}
\end{equation}
Its kernel is $\m_w$, and the inclusion $K_H\subseteq\OO_w$ is a
coefficient-field section. No section of the ordered group quotient
$\Gamma\to\Gamma/H$ is required.

To verify multiplicativity, observe that the support of an element of
$\OO_w$ contains only nonnegative quotient values. A product exponent
can belong to $H$ only if both summand exponents belong to $H$. Thus its
$H$-part is exactly the product of the $H$-parts.

\begin{lemma}[Unrestricted coarse formal evaluation]\label{lem:coarse-eval}
Let $z_1,\ldots,z_d\in\m_w$. For every formal series
\[
 F(Z)=\sum_{\nu\in\N^d}a_\nu Z^\nu\in K_H[[Z_1,\ldots,Z_d]],
\]
the family $(a_\nu z^\nu)_\nu$ is strongly Hahn-summable in $K_\Gamma$.
There is no boundedness assumption on $v(a_\nu)$ in $H$.
Evaluation is a $K_H$-algebra homomorphism. It commutes with formal
composition when the inner series have zero constant coefficient.
\end{lemma}

\begin{proof}
Project the supports of the nonzero $z_i$ into $\Gamma/H$ and let $S$
be their union. A monotone image of a well-ordered set is well ordered:
the least preimage of any nonempty chosen subset maps to its least
image. Since $w(z_i)>0$, every element of $S$ is strictly positive.
Apply \cref{lem:neumann} in the quotient group. The possible quotient
values form the well-ordered monoid $S^*$. At a fixed quotient value
$\bar\gamma$, there are only finitely many decompositions into members
of $S$, including word lengths. In particular, there are only finitely
many multiindices $\nu$ that can contribute there.

For each of these finitely many choices, taking the indicated quotient
fibers of the supports of the $z_i$ and multiplying by $a_\nu$ involves
ordinary finite Hahn products. These are well-defined and have
well-ordered support within the fiber, with finite coefficient
contributions. The finite union over the choices has the same properties.
A subset of an ordered group with well-ordered quotient image and
well-ordered support in every quotient fiber is well ordered: minimize
first the quotient and then the fiber. Hence the union of all supports
is well ordered in $\Gamma$, and each actual exponent receives finitely
many contributions.

For addition and multiplication, use the same argument on the jointly
indexed families. For composition, the zero-constant hypothesis ensures
that at any bounded word length only finitely many formal substitutions
can contribute. The quotient Neumann bound again supplies this word
length bound. Thus the coefficientwise identities survive evaluation.
\end{proof}

\begin{example}[Why coefficient bounds are unnecessary]\label{ex:unbounded}
Let $\Gamma=\Z\oplus_{\mathrm{lex}}\Z$, $H=\{0\}\oplus\Z$,
and $z=t^{(1,0)}$. Then
\[
 \sum_{n\ge0}t^{(0,-n^2)}z^n
       =\sum_{n\ge0}t^{(n,-n^2)}
\]
is strongly summable, although its coefficient valuations
$(0,-n^2)$ are unbounded below in $H$. The first coordinate separates
the degrees. This is precisely the type of separation used in the
uniformization proof.
\end{example}

\subsection{The rank-one field at the \texorpdfstring{$q$}{q}-scale}

Now set $H=H_\alpha$. Use the embedding
$H/H_\alpha^-\hookrightarrow\R$ from \cref{sec:hahn} to define a
real-valued valuation $v_1$ on $K_H$, with $v_1(t^\alpha)=1$.

\begin{lemma}[Completeness of the residue workspace]\label{lem:rank-one}
The field $K_H$, equipped with the absolute value $|a|_1=e^{-v_1(a)}$,
is complete. Every $q_0\in K_H$ with $v(q_0)=\alpha$ has $|q_0|_1<1$.
\end{lemma}

\begin{proof}
Let $(a_n)$ be a Cauchy sequence. Choose a subsequence $(a_{n_k})$ whose
successive differences $d_k=a_{n_{k+1}}-a_{n_k}$ satisfy $v_1(d_k)>k$.
Then $v(d_k)>k\alpha$. The sequence $k\alpha$ is cofinal in $H$:
for $\gamma\in H$, choose an integer bound for $|\gamma|$ and then a
strictly larger integer.

It follows that $(d_k)$ is strongly summable in $K_H$. Indeed, below any
fixed exponent only finitely many of its supports can occur; any
hypothetical descending sequence of support exponents would therefore
lie in a finite union of well-ordered supports. The strong sum
$a_{n_1}+\sum_kd_k$ is the $v_1$-limit of the subsequence. The original
Cauchy sequence has the same limit. The assertion about $q_0$ is
$v_1(q_0)=1$.
\end{proof}

The coefficient field of this rank-one valuation need not be $\C$;
it can itself retain infinitesimal scales. Tate's theorem does not
require it to be $\C$ or require the value group to be discrete.
Its complete-field version is exactly the classical input needed here.

\section{Formal germs and the passage from rank one to arbitrary rank}
\label{sec:germs}

We record carefully how classical analytic identities can be used without
illegally specializing a non-summable family.

\subsection{The classical input, stated with its hypotheses}

\begin{theorem}[Tate; external classical input]\label{thm:classical}
Let $F$ be complete for a nontrivial real-valued non-Archimedean absolute
value, and let $q_0\in F$ satisfy $0<|q_0|<1$. The rational bilateral
series \eqref{eq:X}--\eqref{eq:Y}, interpreted by their convergent tails,
define an analytic group homomorphism
\[
 \phi_{q_0}:F^\times\longrightarrow E_{q_0}(F)
\]
which is onto and has kernel $q_0^\Z$. It extends analytically through
those parameters as a map to the projective curve. The same formulas
are locally analytic in the parameter $q_0$ as well as in $u$.
\end{theorem}

The surjective group statement is Tate's Theorem~1
\cite{Tate}. Local analyticity in both variables follows directly from
the rational tails: on sufficiently small neighborhoods of any non-pole
pair, the finitely many relevant denominators remain nonzero and the
remaining terms tend uniformly to zero geometrically. At a pole one
uses the projective parameter below. This last observation requires no
new uniformization theorem.

The proof of \cref{thm:main} uses \cref{thm:classical} as an input,
not as a theorem claimed to have been independently reproved here.
All additional all-rank summability and lifting steps are supplied.

\subsection{The formal neighborhood of the identity}

\begin{lemma}[A universal formal parameter]\label{lem:formal-identity}
With an external variable $h$, the formal expression
\begin{equation}
 B(q,h)=-\frac{X_q(1+h)}{Y_q(1+h)}\label{eq:B}
\end{equation}
after cancellation of its poles belongs to $h\C[[q,h]]$ and has
linear coefficient $1$ in $h$:
\begin{equation}
 B(q,h)=h+h^2 C(q,h),\qquad C\in\C[[q,h]].\label{eq:Blinear}
\end{equation}
It has a unique inverse under composition in the second variable,
$A(q,z)=z+z^2D(q,z)\in z\C[[q,z]]$.
For any $v(q)>0$ and $v(h)>0$, both series evaluate strongly, and
\begin{equation}
 v(B(q,h))=v(h),\qquad \lc(B(q,h))=\lc(h).\label{eq:isometry}
\end{equation}
The germ $h\mapsto\Phi_q(1+h)$ identifies the multiplicative formal
group with the formal group of $E_q$ at $O$.
\end{lemma}

\begin{proof}
Expand the positive tails with $u=1+h$, treating $q,h$ as formal
variables. They belong to $q\C[[q,h]]$. Hence
\begin{align*}
 h^2X_q(1+h)&=1+h+h^2R_X(q,h),\\
 h^3Y_q(1+h)&=-(1+h)^2+h^3R_Y(q,h),
\end{align*}
where $R_X,R_Y\in q\C[[q,h]]$. The denominator obtained after cancelling
poles has constant coefficient $1$. This proves
\eqref{eq:Blinear}, including the exact coefficient $1$ for every $q$.
Formal coefficient recursion gives the inverse $A$, since this linear
coefficient is a unit.

Positive-support calculus evaluates $B$, $A$, and their composition
identities at the specified arguments. The remainder $h^2 C(q,h)$ has
valuation at least $2v(h)>v(h)$, which proves \eqref{eq:isometry}.

The formal-group identity is a universal identity of formal series in
$q,h_1,h_2$. For ordinary complex $q$ and sufficiently small $h_1,h_2$
it follows from classical multiplicative uniformization. Uniqueness of
Taylor coefficients gives the formal identity; alternatively it is the
standard formal Tate-curve construction \cite{Tate,Conrad}. Strong
evaluation of the jointly indexed positive-support family proves it in
$K_\Gamma$.
\end{proof}

Here is a direct explanation of why this formal germ parametrizes all
points in a coarse identity fiber, rather than just some of them.
Use projective coordinates $[X:Y:Z]$ near $O=[0:1:0]$ and set
\[
 z=-X/Y=-x/y,\qquad r=-Z/Y=-1/y.
\]
The curve equation becomes
\begin{equation}
 r(1-z)=z^3+a_4zr^2+a_6r^3.\label{eq:localcurve}
\end{equation}
It has a unique solution $r\in z^3\C[[q,z]]$ by coefficient recursion.
For $z\in\m_w$, positive-support evaluation gives such a solution.
Uniqueness also holds among actual $r\in\m_w$: subtracting two
solutions factors their difference times
\[
 1-z-a_4z(r_1+r_2)-a_6(r_1^2+r_1r_2+r_2^2),
\]
which is a $w$-unit because its residue is $1$. Thus two points reducing
to $O$ and having the same $z$ coincide. Evaluation of the inverse
$A(q,z)$ supplies the unique corresponding $h\in\m_w$.

\subsection{Taylor prolongation through coarse infinitesimals}

Let $w=v\bmod H_\alpha$, set $F=K_{H_\alpha}$, and put
$q_0=\res_{H_\alpha}(q)$. Then
\[
 q=q_0(1+\rho),\quad \rho\in\m_w,
 \qquad
 u=u_0(1+\sigma),\quad \sigma\in\m_w,
 \quad u_0=\res_{H_\alpha}(u)\in F^\times
\]
for every $u\in U_q$.

\begin{lemma}[Agreement of Taylor and Hahn evaluation]\label{lem:agreement}
The classical local Taylor germ of the projective map
$(Q,U)\mapsto\phi_Q(U)$ at $(q_0,u_0)$, evaluated at
$(\rho,\sigma)$ by \cref{lem:coarse-eval}, is well defined and equals
$\Phi_q(u)$ as defined by the literal Hahn formulas.
This remains true when $u_0\in q_0^\Z$, using a projective local chart.
\end{lemma}

\begin{proof}
First suppose $u_0\notin q_0^\Z$. Replace $(u,u_0)$ by
$(u q^m,u_0 q_0^m)$ so that
$0\le v(u_0)<\alpha$. This is a finite reindexing of the bilateral
formulas, and its relative error is still in $\m_w$.

The isolated rational terms have nonzero constant denominators in $F$.
Their Taylor series in $\rho,\sigma$ evaluate to the actual rational
terms by \cref{lem:coarse-eval}. For the $n$th positive rational tail,
expand first in the relative variables. At every fixed Taylor
multiindex, the coefficient lies in $F$ and its fine valuation tends
to $+\infty$ cofinally in $H_\alpha$ as $n\to\infty$.
More explicitly, after discarding finitely many $n$, the denominators
$1-q_0^nu_0$ and $1-q_0^n/u_0$ are fine units, and a coefficient of
any fixed relative Taylor degree has valuation at least
$n\alpha-v(u_0)$, or a larger positive multiple for the $g$-tail.
Differentiating relative powers only contributes ordinary integer
coefficients and additional factors of these small rational arguments.
The finitely many discarded denominators affect no tail conclusion.

It follows that the sum over $n$ defining every Taylor coefficient is a
strong sum in $F$ and agrees with its rank-one analytic sum, by
\cref{lem:rank-one}. Next fix a quotient exponent for $w$. The quotient
Neumann bound in \cref{lem:coarse-eval} permits only finitely many
Taylor multiindices and quotient-support decompositions at this value.
For each such choice, the remaining sum over $n$ is the cofinally
bounded sum just described. Thus the family jointly indexed by Taylor
degree, tail index, and support decompositions is strongly summable.
Fubini regrouping for this joint family proves that Taylor evaluation
and the original Hahn-tail evaluation coincide. The $s_1$ family has
the same, simpler estimates.

If $u_0=q_0^m$, divide $u$ by $q^m$. The result is $1+h$ with
$h\in\m_w$. Apply \cref{lem:formal-identity} and the regular projective
chart \eqref{eq:localcurve}. The poles of $X,Y$ disappear in this chart,
and the same formal expressions are the Taylor germs of the classical
projective map. This proves the claim, including $h=0$.
\end{proof}

\begin{remark}[The interchange that must not be omitted]
Formal coefficients in $F$ can have arbitrarily negative fine
valuations, as \cref{ex:unbounded} shows. Ordinary positive-valuation
substitution over $\C$ alone does not justify their evaluation. Also,
a sequence cofinal in $H_\alpha$ need not be cofinal in $\Gamma$.
The proof above uses two separate finiteness statements: quotient
Neumann finiteness controls the coarse Taylor degrees, while cofinality
inside $H_\alpha$ controls the rational tail index. Neither statement
substitutes for the other.
\end{remark}

\section{Proof of arbitrary-rank uniformization}\label{sec:uniformization}

\subsection{The curve equation and the group law}

Take $F=K_{H_\alpha}$, with its complete rank-one coarsening, as in
\cref{lem:rank-one}. Since $v(q_0)=\alpha$, the classical theorem
applies to $q_0$. The projective Taylor germs in
\cref{lem:agreement} satisfy the curve equation: it is an identity of
classical analytic germs, hence an identity of their Taylor series.
Evaluation preserves these polynomial identities. Thus $\Phi_q$ maps
into $E_q(K)$.

The same argument proves
\begin{equation}
 \Phi_q(uv)=\Phi_q(u)+\Phi_q(v)\qquad(u,v\in U_q).\label{eq:grouplaw}
\end{equation}
Here one must use regular local charts on the elliptic curve, not the
chord-slope formula at a vanishing denominator. The group operation on a
smooth elliptic curve is a regular morphism. At the residue points
$\phi_{q_0}(u_0)$, $\phi_{q_0}(v_0)$ and their sum, choose affine
neighborhoods for which it is given by rational expressions with
nonzero residue denominators. Its local Taylor identity with
$\phi_Q(UV)$ follows from \cref{thm:classical}. All denominators remain
$w$-units after coarse evaluation, and \cref{lem:coarse-eval} preserves
the identity. At the identity, use \eqref{eq:localcurve}. This includes
$u_0v_0\in q_0^\Z$, opposite points, equal points, and every other
case where a particular affine addition formula would be undefined.
By \cref{lem:agreement}, these evaluated germs are the Hahn-defined
map, proving \eqref{eq:grouplaw} everywhere on $U_q$.

For clarity, this is not an appeal to a new ``identity principle on
all surcomplex numbers.'' It is an ordinary identity of rank-one
analytic germs, followed by an explicitly proved algebra homomorphism
on their formal Taylor rings.

\subsection{Good reduction for the coarse valuation}

By \eqref{eq:Delta}, $v(\Delta(q))=\alpha\in H_\alpha$, so
$w(\Delta(q))=0$. All coefficients of the projective Weierstrass equation
belong to $\OO_w$, and its discriminant is a unit there. It is therefore
a smooth proper elliptic curve over $\OO_w$, with residue curve
$E_{q_0}$ over $F$.

Every $K$-point has a well-defined reduction to that curve. In the
present projective model, one can see existence directly: scale its
three homogeneous coordinates so that their minimum $w$-value is
zero, and reduce the resulting integral coordinates. The equation
survives reduction, and at least one reduced coordinate is nonzero.
Smooth properness also implies that reduction respects the group law.
We write
\[
 \mathrm{red}:E_q(K)\longrightarrow E_{q_0}(F),
 \qquad E_1(K)=\ker(\mathrm{red}).
\]
The reduction compatibility supplied by \cref{lem:agreement} is
\begin{equation}
 \mathrm{red}(\Phi_q(u))=\phi_{q_0}(u_0).\label{eq:redcompat}
\end{equation}

\begin{proposition}[The whole coarse identity fiber]\label{prop:identityfiber}
The homomorphism $h\mapsto\Phi_q(1+h)$ induces an isomorphism
\begin{equation}
 1+\m_w\ \simeq\ E_1(K).\label{eq:identityfiber}
\end{equation}
For the local coordinate $z=-x/y$, it has the explicit inverse
$h=A(q,z)$ and satisfies $v(z)=v(h)$ for $h\ne0$.
\end{proposition}

\begin{proof}
A point reducing to $O=[0:1:0]$ lies in the chart of
\eqref{eq:localcurve}, with $z,r\in\m_w$. The discussion following
\cref{lem:formal-identity} proves that $z$ uniquely determines the
point and that the inverse formal parameter supplies a unique
$h\in\m_w$. Formal composition identities show that these maps are
mutual inverses. The formal-group identity gives the group law; the
valuation assertion is \eqref{eq:isometry}.
\end{proof}

\subsection{Surjectivity}

Let $P\in E_q(K)$ be arbitrary, and reduce it to $P_0\in E_{q_0}(F)$.
By the classical Tate theorem there is $u_0\in F^\times$ with
$\phi_{q_0}(u_0)=P_0$. Regard $u_0$ as an element of $K$. Its valuation
belongs to $H_\alpha$, so $\Phi_q(u_0)$ is defined. By
\eqref{eq:redcompat} the point
\[
 P-\Phi_q(u_0)
\]
belongs to $E_1(K)$. By \cref{prop:identityfiber} it equals
$\Phi_q(1+h)$ for a unique $h\in\m_w$. The group law now gives
\[
 P=\Phi_q\bigl(u_0(1+h)\bigr).
\]
The product parameter is in $U_q$. This proves surjectivity, including
points with coordinates outside every scale bounded by powers of $q$.

When $H_\alpha=\Gamma$, the valuation $w$ is trivial and
$\m_w=0$. The proof reduces to the classical theorem on $K_H=K$
with its rank-one coarsening. No separate nontrivial-coarse-valuation
hypothesis was smuggled into the argument.

\subsection{The kernel}

Suppose $\Phi_q(u)=O$. Reducing and using the classical kernel shows
$u_0=q_0^m$ for an integer $m$. Therefore $u/q^m\in1+\m_w$.
Reindexing already gives $\Phi_q(q^m)=O$; the group law and
\cref{prop:identityfiber} imply $u/q^m=1$. Conversely every $q^m$
is mapped to $O$ by definition. Thus $\ker\Phi_q=q^\Z$.
Together with the support theorems and surjectivity, this completes
the proof of \cref{thm:main}.

\subsection{The exact reduction diagram}

The proof gives more than an abstract isomorphism.

\begin{corollary}[Reduction and infinitesimal lifting]\label{cor:diagram}
There is a commutative diagram of exact sequences
\[
\begin{tikzcd}[column sep=small,row sep=large]
1\arrow[r]&1+\m_w\arrow[r]\arrow[d,"\sim"']&U_q/q^\Z
 \arrow[r,"\res_H"]\arrow[d,"\sim"']&F^\times/q_0^\Z
 \arrow[r]\arrow[d,"\sim"]&1\\
0\arrow[r]&E_1(K)\arrow[r]&E_q(K)\arrow[r,"\mathrm{red}"]&E_{q_0}(F)
 \arrow[r]&0 .
\end{tikzcd}
\]
The left-hand isomorphism is given by a support-controlled universal
formal series and preserves the original valuation of the local
identity parameter.
\end{corollary}

\begin{proof}
Surjectivity of the top residue map follows by taking representatives
in $F$. Its kernel consists of those classes whose residue is a power
of $q_0$; division by the corresponding power of $q$ gives a unique
representative in $1+\m_w$. Uniqueness uses $v(q_0)>0$.
The lower row and all vertical maps were established above.
\end{proof}

\section{Scalar extension and the valuation-circle quotient}\label{sec:extension}

\subsection{Strong functoriality}

\begin{theorem}[Exponent embeddings preserve uniformization]\label{thm:extension}
Let $\iota:\Gamma\hookrightarrow\Gamma'$ be an ordered group embedding.
The coefficient-preserving map
\[
 \iota_*:K_\Gamma\hookrightarrow K_{\Gamma'},\qquad
 \sum a_\gamma t^\gamma\longmapsto\sum a_\gamma t^{\iota(\gamma)}
\]
is a field embedding preserving strong sums. If $q'=\iota_*(q)$, then
\begin{equation}
 \iota_*(\Phi_q(u))=\Phi_{q'}(\iota_*(u))\qquad(u\in U_q),
 \label{eq:functorial}
\end{equation}
with the corresponding identities for $\Theta,X,Y$ where defined.
Moreover
\begin{equation}
 \iota^{-1}(H_{\iota(\alpha)})=H_\alpha.\label{eq:Hpullback}
\end{equation}
In the larger field the full admissible domain is
$v^{-1}(H_{\iota(\alpha)})$, not in general all of $K_{\Gamma'}^\times$.
\end{theorem}

\begin{proof}
An order embedding preserves and reflects well-ordering and preserves
finite coefficient contributions. It therefore preserves the defining
Hahn operations and every strongly summable family. Apply this to the
bilateral formulas, or to their positive-tail forms. The point at
infinity and the powers of $q$ are also preserved. Finally, inequalities
$|\iota(\gamma)|\le N\iota(\alpha)$ are equivalent to
$|\gamma|\le N\alpha$, proving \eqref{eq:Hpullback}. The exact-domain
theorems over $\Gamma'$ give the last assertion.
\end{proof}

This theorem has a useful asymmetric meaning. An old admissible
parameter never becomes inadmissible under an ordered exponent embedding.
But the assertion ``every nonzero parameter is admissible'' may fail
after enlargement: new valuations can lie outside the old $q$-scale.
The curve acquires new points, yet its correct, enlarged admissible
subgroup still reaches all of them.

\begin{example}[A rank-one base and a rank-two extension]\label{ex:extension}
Embed $\R$ as the second summand in
$\Gamma'=\R\oplus_{\mathrm{lex}}\R$, and take $q=t^{(0,1)}$.
Over $\C((t^\R))$, every nonzero parameter is admitted. Over
$K_{\Gamma'}$, a parameter of valuation $(1,0)$ is not admitted.
A parameter $1+t^{(1,0)}$ is admitted and yields a new point reducing
to $O$ for the coarse first-coordinate valuation. Thus excluding a new
multiplicative scale does not exclude its infinitesimal displacement
from a nonzero old parameter.
\end{example}

\subsection{A canonical circularly ordered value quotient}

Let $\OO_v^\times=\{u\in K^\times:v(u)=0\}$ be the units of the
original Hahn valuation ring.

\begin{theorem}[The valuation quotient]\label{thm:tropical}
There is a canonical surjective group homomorphism
\begin{equation}
 \trop_q:E_q(K)\longrightarrow H_\alpha/\Z\alpha,
 \qquad \trop_q(\Phi_q(u))=v(u)+\Z\alpha,\label{eq:tropical}
\end{equation}
and an exact sequence
\begin{equation}
 1\longrightarrow\OO_v^\times\xrightarrow{\ \Phi_q\ }
 E_q(K)\xrightarrow{\ \trop_q\ }H_\alpha/\Z\alpha
 \longrightarrow0.\label{eq:tropical-seq}
\end{equation}
Every value class has a unique representative in $[0,\alpha)$.
The ordering of that interval defines a translation-invariant circular
order on the quotient, with addition followed by reduction modulo
$\alpha$.
\end{theorem}

\begin{proof}
Two parameters represent the same point exactly when their quotient is
a power of $q$, so their valuations differ by an integer multiple of
$\alpha$. This proves well-definedness. Every $\gamma\in H_\alpha$
is the valuation of $t^\gamma$, proving surjectivity.
If $v(u)=m\alpha$, the same point is represented by $u/q^m$, which is
a valuation-zero unit. Conversely units map to zero in the value
quotient. Their map into $E_q(K)$ is injective, since a power of $q$
has valuation zero only when it is $1$. Finite normalization gives the
unique representative. The usual cyclic order of three representatives
in a fundamental half-open interval is independent of simultaneous
translation, by cutting the interval at the translating point.
\end{proof}

There is a second exact sequence of additive groups
\begin{equation}
 0\longrightarrow H_\alpha^-\longrightarrow H_\alpha/\Z\alpha
 \longrightarrow
 (H_\alpha/H_\alpha^-)/\Z\bar\alpha\longrightarrow0.
 \label{eq:circlelayers}
\end{equation}
Indeed $H_\alpha^-\cap\Z\alpha=0$, and the kernel computation is
immediate. This separates infinitesimal value directions from the
rank-one circular quotient.

When the relevant rank-one value group is all of $\R$, this value
quotient is the familiar real circle of length $v_1(q)$. Over a complete
algebraically closed rank-one field it agrees with the value-circle
description of the Tate curve skeleton \cite{BPR}.
At arbitrary rank, \eqref{eq:tropical-seq} is an \emph{algebraic and
circular-order} assertion. It does not claim the existence of an ordinary
real-metric Berkovich skeleton with all of $H_\alpha$ as its distance
set. Higher-rank analytifications require their own definitions
\cite{FR}.

\section{Recovering the Tate parameter without adding scales}\label{sec:j}

\subsection{The inverse modular series}

Write $J=1/j$. The formal series
\[
 J(q)=\frac{\Delta(q)}{c_4(q)^3}\in q\Z[[q]]
\]
has linear coefficient $1$, hence a unique compositional inverse
$Q(J)\in J\Z[[J]]$. Its first terms are
\begin{equation}
 Q(J)=J+744J^2+750420J^3+872769632J^4
       +1102652742882J^5+\cdots.\label{eq:qJ}
\end{equation}
The notation $Q$ here denotes a formal series, not the rational field.

\begin{theorem}[Unique, support-preserving recovery]\label{thm:recover}
Given $j\in K_\Gamma$ with $v(j)<0$, there is a unique $q\in K_\Gamma$
with $v(q)>0$ and $j(q)=j$. It is $q=Q(j^{-1})$ and satisfies
$v(q)=-v(j)$. If $G(a)$ denotes the additive subgroup of $\Gamma$
generated by $\supp(a)$, then
\begin{equation}
 G(q)=G(j).\label{eq:supportgroup}
\end{equation}
More precisely,
\begin{equation}
 \supp(q)\subseteq
 \bigl(\supp(j^{-1})\bigr)^*\setminus\{0\}.\label{eq:qmonoid}
\end{equation}
\end{theorem}

\begin{proof}
Since $v(j^{-1})>0$, positive-support calculus evaluates $Q$ and both
formal inverse identities. Thus $J(q)=j^{-1}$ and $j(q)=j$.
The leading coefficient of $Q$ is $1$, so $v(q)=v(j^{-1})$.
If another positive-valuation $q'$ gives $j(q')=j$, evaluating the
identity $Q(J(q'))=q'$ proves equality. This argument does not choose
a branch of a transcendental function.

Formula \eqref{eq:qmonoid} follows directly from the positive powers of
$j^{-1}$. Inversion of a Hahn series uses only the additive subgroup of
its support, so $G(j^{-1})\subseteq G(j)$. Therefore $G(q)\subseteq G(j)$.
Conversely $j(q)=q^{-1}$ times a power series in $q$ with nonzero
constant coefficient. Inversion and strong positive evaluation use only
$G(q)$, giving $G(j)\subseteq G(q)$.
\end{proof}

The inverse modular expansion is classical. The support equality is the
useful extra statement in a symbolic or nondivisible Hahn workspace:
one need not pass to the divisible hull merely to recover the elliptic
period parameter.

\subsection{Which elliptic curves this covers}

\begin{corollary}[Negative-$j$ uniformization]\label{cor:negativej}
Assume $\Gamma$ is divisible. Every elliptic curve $E/K_\Gamma$ with
$v(j(E))<0$ is $K_\Gamma$-isomorphic to $E_q$, for the unique parameter
of \cref{thm:recover}. After any such choice of isomorphism,
$E(K_\Gamma)$ is parametrized by $U_q/q^\Z$.
\end{corollary}

\begin{proof}
A Hahn field with algebraically closed coefficient field of
characteristic zero and divisible exponent group is algebraically
closed \cite{Poonen}. Elliptic curves over an algebraically closed field
of characteristic zero with the same $j$-invariant are isomorphic
\cite{SilvermanArithmetic}. Apply this to $E$ and $E_q$.
\end{proof}

The parameter $q$ is canonical, but an isomorphism $E\simeq E_q$ need
not be canonical. The hypothesis on $j$ excludes $j=0$ and $1728$, but
does not remove the usual sign ambiguity of an elliptic isomorphism.
For nondivisible $\Gamma$, \cref{thm:recover} still holds, while a
curve with the same $j$ can be a nontrivial twist; we do not silently
identify all such twists with the split Tate curve over the original
field. Curves with $v(j)\ge0$ are outside this positive-$q$ theorem.

\section{Torsion fields and exact exponent ramification}\label{sec:torsion}

This section retains nondivisible exponent groups deliberately. Replacing
$\Gamma$ by its divisible hull at the start would make every degree
computed here equal to one after base extension.

\subsection{Roots of a Hahn parameter}

\begin{lemma}[A single radical extension]\label{lem:radical}
Let $n\ge2$, $q\in K_\Gamma$ with $v(q)=\alpha>0$, and
$\Lambda=\Gamma+\Z\alpha/n$. Then for any choice of $q^{1/n}$ in
$K_{\Gamma\otimes\Q}$,
\begin{equation}
 K_\Gamma(q^{1/n})=K_\Lambda,\qquad
 [K_\Lambda:K_\Gamma]=[\Lambda:\Gamma]
       =\ord(\alpha+n\Gamma).\label{eq:radical}
\end{equation}
The extension is cyclic, with Galois group
$\operatorname{Hom}(\Lambda/\Gamma,\C^\times)$; it has unchanged
residue field $\C$ and value group $\Lambda$.
\end{lemma}

\begin{proof}
Write $q=c t^\alpha(1+\varepsilon)$, $c\in\C^\times$, $v(\varepsilon)>0$.
Choose $c^{1/n}\in\C$. The binomial series gives
$b=(1+\varepsilon)^{1/n}\in K_\Gamma$, with $b^n=1+\varepsilon$.
Consequently $q^{1/n}=c^{1/n}t^{\alpha/n}b$ up to a root of unity
already in $\C$, so adjoining $q^{1/n}$ is the same as adjoining
$t^{\alpha/n}$.

The group $\Lambda/\Gamma$ is finite cyclic. If its order is $d$,
partition any Hahn support in $\Lambda$ among its $d$ cosets. Factoring
a fixed monomial from each part gives a finite sum of $d$ monomials
with coefficients in $K_\Gamma$. These monomials are linearly independent
over $K_\Gamma$, since their supports lie in disjoint cosets. Thus the
entire Hahn field $K_\Lambda$, not just a smaller subfield with finite
supports, is a degree-$d$ extension generated by $t^{\alpha/n}$.

Multiplying the monomial of exponent $\lambda$ by a character of
$\lambda+\Gamma$ defines a field automorphism, preserving strong sums.
There are $d$ such characters because $\C$ contains all roots of unity,
and they give all automorphisms. The coefficient residue field remains
$\C$ and the ramification index is $d$, equal to the field degree.
Finally $d$ is the least positive integer with $d\alpha\in n\Gamma$,
which is the displayed order in $\Gamma/n\Gamma$.
\end{proof}

\subsection{A primitive torsion point}

\begin{theorem}[The complete torsion field]\label{thm:torsion}
Under the hypotheses of \cref{lem:radical}, the $n^2$ points
\begin{equation}
 \Phi_q(\zeta_n^a q^{b/n}),\qquad 0\le a,b<n,\label{eq:torsionpts}
\end{equation}
are precisely $E_q[n]$. Their field of definition is $K_\Lambda$,
and already
\begin{equation}
 K_\Gamma\bigl(\Phi_q(q^{1/n})\bigr)=K_\Lambda.\label{eq:primitivepoint}
\end{equation}
Here adjoining a point means adjoining its affine coordinates.
\end{theorem}

\begin{proof}
Every parameter in \eqref{eq:torsionpts} lies in the admissible group of
$K_\Lambda$, since its valuation is $b\alpha/n$. Its $n$th power is
$q^b$, so the group law makes its image $n$-torsion. If two of the
points coincide, injectivity modulo $q^\Z$ gives
\[
 \zeta_n^{a-a'}q^{(b-b')/n}=q^m
\]
for some integer $m$. Comparing valuations gives $b-b'=nm$; the chosen
ranges force $m=0$ and $b=b'$, and then $a=a'$. Thus there are $n^2$
distinct points. In characteristic zero an elliptic curve has exactly
$n^2$ geometric $n$-torsion points, proving completeness of the list.
They all have coordinates in $K_\Lambda$ by strong evaluation.

To prove the stronger primitive assertion, let
$L=K_\Gamma(q^{1/n})=K_\Lambda$ and let $\sigma\in\Gal(L/K_\Gamma)$
fix the point $P=\Phi_q(q^{1/n})$. The character description in
\cref{lem:radical} implies that $\sigma$ preserves every strong sum,
so it commutes with the coordinate formulas. Therefore
\[
 \Phi_q(\sigma(q^{1/n}))=\Phi_q(q^{1/n}).
\]
The quotient $\sigma(q^{1/n})/q^{1/n}$ is a root of unity. By the
uniformization kernel it must also be a power of $q$. A root of unity
has valuation zero, whereas $v(q^m)=m\alpha$; hence that power is $1$.
Thus $\sigma$ fixes $q^{1/n}$ and is the identity. The stabilizer of
$P$ is trivial, so Galois theory gives \eqref{eq:primitivepoint}.
The full torsion field contains this primitive point field and is
contained in $L$, proving equality.
\end{proof}

\begin{corollary}[A finite-lattice degree formula]\label{cor:lattice}
If $\Gamma\cong\Z^r$ as an abelian group, with any compatible total
order, and $\alpha$ has coordinates $(a_1,\dots,a_r)$, then
\begin{equation}
 [K_\Gamma(E_q[n]):K_\Gamma]
      =\frac{n}{\gcd(n,a_1,\dots,a_r)}.\label{eq:gcddegree}
\end{equation}
\end{corollary}

\begin{proof}
The order of $(a_1,\ldots,a_r)$ in $(\Z/n\Z)^r$ is the least $d>0$
for which $n$ divides every $da_i$. This is the displayed integer.
\end{proof}

For example, $n=12$ and $\alpha=(6,4)$ give degree $6$; $n=4$ and
$\alpha=(0,6)$ give degree $2$ in the lexicographic group $\Z^2$.
The entire unit part of $q$ can have infinite support and does not
alter these degrees. It is removed by a binomial root already in the
base field. Over $K_\Q$ every $n$th root scale is already available,
so the corresponding degrees are one. This explains why scalar
extension can erase arithmetic information while preserving all
uniformization identities.

\section{Transport to surreal and surcomplex numbers}\label{sec:surcomplex}

\subsection{Set-sized normal-form workspaces}

The Conway normal form writes a surreal as a reverse-well-ordered sum
of real multiples of monomials $\omega^\gamma$. With the convention
$t^\gamma=\omega^{-\gamma}$, increasing well-ordered Hahn supports
correspond to these normal forms. Combining the real and imaginary
normal forms produces complex coefficients and gives the surcomplex
version \cite{Gonshor,BournezGuilmant}.

If a set of surcomplex parameters and point coordinates is specified,
the union of all of its normal-form exponent supports is a set. Its
additive subgroup $\Gamma\subseteq(\No,+)$ is a set as well. The
associated normal-form map embeds $K_\Gamma$ into $\No[i]$ and contains
all the specified elements. Enlarge to the divisible hull only when
algebraic closedness is actually needed.

The proof of \cref{thm:main} works in this set-sized workspace. If a
larger workspace is chosen, \cref{thm:extension} shows that the formulas
and their values agree. There is therefore no choice-dependent global
summation rule and no need to regard a proper class as a set.

\begin{theorem}[Surcomplex uniformization]\label{thm:surcomplex}
For every $q\in\No[i]^\times$ with $\alpha=v(q)>0$, the literal Hahn
Tate map has exactly the admissible domain in \eqref{eq:No-main} and
induces the group isomorphism stated there. This domain is always a
proper subclass of $\No[i]^\times$.
For every elliptic curve $E/\No[i]$ with $v(j(E))<0$, the unique
parameter $q=Q(j(E)^{-1})$ yields an isomorphic Tate curve.
\end{theorem}

\begin{proof}
The domain test for one parameter involves only $\alpha$, its valuation,
and ordinary integer inequalities. Its truth is unchanged in any
workspace containing those data. The same holds for all the strong
families in the formulas. For surjectivity, put $q$ and a given point's
coordinates in one workspace and apply \cref{thm:main} there. For the
kernel and the group law, place the finitely many relevant parameters
in a common workspace. These are instances of \cref{thm:extension}.

To see properness, take a positive surreal value
$\gamma=\omega\alpha$. It exceeds $N\alpha$ for every ordinary $N$.
The monomial $u=\omega^{-\gamma}$ lies outside the domain, by the exact
support theorem.

For the final assertion put the Weierstrass coefficients of $E$ in a
workspace and take its divisible hull. The algebraically closed Hahn
field there contains an isomorphism supplied by
\cref{cor:negativej}. Its image in $\No[i]$ proves the assertion.
The parameter obtained by the inverse series is independent of the
workspace, again by strong functoriality.
\end{proof}

\subsection{A forbidden monomial and an admitted fine displacement}

\begin{example}[Two ways to use $\omega^{-\omega}$]\label{ex:omega}
Let $q=\omega^{-1}$, so $\alpha=1$, and put
$\varepsilon=\omega^{-\omega}$. The parameter $u=\varepsilon$ is
inadmissible: $v(u)=\omega>N$ for all ordinary $N$.
For the theta family the exponents along $n=-k$ are
\[
 \frac{k(k+1)}2-k\omega,
\]
a strictly decreasing sequence. For the $X$ family they include
$\omega-k$, again strictly decreasing. No cancellation argument can
make either original family strongly summable.

In contrast $u=1+\varepsilon$ has valuation zero and is admissible.
The point has
\begin{align*}
 X_q(1+\varepsilon)&=\varepsilon^{-2}+\varepsilon^{-1}
                         +\text{a positive-valuation tail},\\
 Y_q(1+\varepsilon)&=-\varepsilon^{-3}-2\varepsilon^{-2}
                         -\varepsilon^{-1}
                         +\text{a positive-valuation tail},\\
 -X_q(1+\varepsilon)/Y_q(1+\varepsilon)
                    &=\varepsilon+O(\varepsilon^2).
\end{align*}
The last $O$ means valuation at least $2\omega$, not an ordinary
real asymptotic estimate. Thus an infinitesimal displacement far finer
than the $q$-scale is retained exactly, and its image has correspondingly
large elliptic coordinates.
\end{example}

\subsection{Why fine-topological convergence is not being used}

For $q=\omega^{-1}$, the powers $q^n$ tend to zero in a rank-one
workspace such as $\C((t^\R))$. They do not tend to zero when neighborhoods
can test against every surreal positive scale: none of the ordinary
$q^n$ is smaller than $\omega^{-\omega}$. The repository already
emphasizes this distinction between its fixed rank-one analysis and
full fine topology \cite{RepoRankOne}.

The present proof uses ordinary convergence only in the named complete
rank-one coarsening of $K_{H_\alpha}$. Everywhere else it uses strong
Hahn summation or finite algebra. Its finite-tail bound is a statement
in the ordered value group; it is not relabeled as an all-scale limit.
Neither a global surcomplex exponential nor a value of $\sin\omega$
is required. Multiplicative periods $q^\Z$ replace those additive
analytic constructions.

\section{Exact computation, proof dependencies, and implementation}
\label{sec:computation}

\subsection{Finite checks included with this manuscript}

The accompanying program \texttt{code/check\_identities.py} uses only
exact rational symbolic arithmetic. Its recorded output is in
\texttt{data/verification.json} and \texttt{data/verification.log}.
The run used Python~3.13.5 and SymPy~1.14.0; the JSON records the actual
runtime version and is authoritative if the script is rerun elsewhere.
The checks do not require a numerical value of $q$ or a realization of
an arbitrary ordered group.

For clarity, the coordinate coefficients used by the script are
\begin{align}
 [q^d]X_q(u)&=\sum_{m\mid d}m(u^m+u^{-m}-2),\quad d\ge1,
 \label{eq:Xcoeff}\\
 [q^d]Y_q(u)&=\sum_{m\mid d}
 \left(\frac{m(m-1)}2u^m-\frac{m(m+1)}2u^{-m}+m\right),\quad d\ge1.
 \label{eq:Ycoeff}
\end{align}
The degree-zero coefficients are $f(u)$ and $g(u)$. These identities
come directly from the positive rational tails; they are also recorded
in the contemporary Tate formulas of \cite{EV}. The script checks the
curve equation, \eqref{eq:diff}, and the inversion identity through
$q^5$ as rational functions of the indeterminate $u$. It checks the
discriminant/product identity through $q^9$, computes $j$ through $q^5$,
and checks inverse modular composition through $J^6$.

The first verified coefficient lists are
\begin{align*}
 a_4(q)&=-5q-45q^2-140q^3-365q^4-630q^5+O(q^6),\\
 a_6(q)&=-q-23q^2-154q^3-647q^4-1876q^5+O(q^6),\\
 \Delta(q)&=q-24q^2+252q^3-1472q^4+4830q^5+O(q^6).
\end{align*}
The formal-parameter sign is checked at $q=0$:
\[
 -\frac{X_0(1+h)}{Y_0(1+h)}=\frac{h}{1+h}
                =h-h^2+h^3-h^4+\cdots.
\]
The test suite also checks finite index-shift instances of the theta
functional equation, explicit decreasing exponent witnesses in a
rank-two lexicographic group, and sample torsion-degree calculations.

These checks are normalization and transcription guards. They do not
prove strong summability, arbitrary-rank surjectivity, an identity at
all orders, or a proper-class transport theorem. Those conclusions
depend on the mathematical arguments above.

\subsection{What an exact symbolic implementation must know}

To represent a concrete instance, a symbolic system needs finite
expressions for the input Hahn series or a reliable lazy coefficient
interface, exact leading exponents and coefficients, and a specified
ordered exponent group. The domain test additionally requires deciding
whether $|\beta|\le N\alpha$ for some ordinary $N$. Comparisons with
individual $N\alpha$ are not, by themselves, a terminating decision
procedure for that existential statement.

For a lexicographically ordered free finite-rank group, the test is
finite. Inspect the first nonzero coordinate of $\alpha$. Values whose
first nonzero coordinate occurs earlier are outside $H_\alpha$;
values supported at that coordinate or later are inside it. In rational
coordinates the finite normalization in \cref{lem:normalize} is then
performed by ordinary floor arithmetic at that coordinate, with the
lower coordinates determining boundary cases. For a general group,
the needed principal-convex-subgroup operation must be part of the
representation contract.

For normalized admissible $u$, an evaluator can return the two isolated
rational terms, a lazy positive-tail object, the support certificate
\eqref{eq:certificate}, and the tail bound \eqref{eq:tail}. That is an
exact expression with an explicit precision policy. A finite truncation
is not an exact elliptic point unless the discarded tail is retained
symbolically. Testing the curve equation numerically after truncation
is not a replacement for the exact identity.

The torsion-degree calculation requires membership in $n\Gamma$ and
finite cyclic order in $\Gamma/n\Gamma$, which is extra group data.
When $\Gamma\cong\Z^r$, \eqref{eq:gcddegree} supplies a direct algorithm.
In a black-box ordered group, it need not be decidable from order
comparisons alone.

\subsection{A dependency-aware formalization plan}

The repository already has a substantial Hahn-support layer, including
positive-support evaluation and regrouping \cite{Repo}. It should be
reused. The new proof obligations naturally separate as follows.

\begin{center}
\begin{tabular}{p{.29\textwidth}p{.62\textwidth}}
\toprule
Layer & Precise obligations \\
\midrule
Ordered groups & Principal convex subgroup, $H_\alpha^-$, finite annulus
normalization, exact theta-leading-exponent obstruction.\\[4pt]
Hahn calculus & Strong coordinate tails, exponent embedding functoriality,
coarse coefficient field, unrestricted coarse formal evaluation.\\[4pt]
Formal algebra & Tate coefficients, local parameter $B$ and its inverse,
projective identity chart, formal-group identities.\\[4pt]
Rank-one bridge & Completeness of $K_{H_\alpha}$ under its rank-one
coarsening and the full classical Tate theorem with explicit hypotheses.\\[4pt]
All-rank geometry & Taylor/Hahn agreement, group law in regular charts,
coarse reduction diagram, surjectivity and exact kernel.\\[4pt]
Arithmetic & Finite coset decomposition of a Hahn extension, character
Galois action, primitive torsion-point stabilizer.\\
\bottomrule
\end{tabular}
\end{center}

This is a plan, not a supplied Lean formalization. In particular, proving
the support lemmas while postulating classical Tate uniformization would
only establish a theorem conditional on that input. A full formal audit
must distinguish that situation from a checked construction of the
rank-one theorem. No unchecked code skeleton is included as though it
were a machine-verified proof.

\section{Limitations and further mathematical questions}\label{sec:limits}

The present results settle the internal domain and surjectivity questions
posed in \cref{sec:intro}. They leave several broader tasks genuinely
separate.

\paragraph{Beyond negative valuation of $j$.}
The Tate parameter $v(q)>0$ produces $v(j)<0$. Curves with nonnegative
valuation of $j$ require a different treatment, such as formal groups
around a good-reduction model and a residue elliptic curve. The
surcomplex field's algebraic closedness does not by itself provide a
single useful transcendental parametrization of that different regime.

\paragraph{Higher-dimensional totally degenerate abelian varieties.}
A natural extension replaces $q^\Z$ by a finite-rank multiplicative
period lattice and the quadratic theta exponent by a positive
polarization form. One would need an exact ordered-group condition for
all lattice directions and a replacement for the single principal
convex subgroup $H_\alpha$. The support proof in this paper suggests
the right obstruction, but it is not a proof of a general multivariable
uniformization theorem.

\paragraph{Intrinsic higher-rank analytic geometry.}
The algebraic quotient $H_\alpha/\Z\alpha$ retains infinitesimal value
directions discarded by its Archimedean quotient. Relating it precisely
to Hahn analytification, adic spaces, or a suitable skeleton is a
separate geometric theorem. The exact sequence \eqref{eq:circlelayers}
is a starting invariant, not a substitute for that theorem.

\paragraph{Functions outside the literal Hahn domain.}
We proved that the displayed bilateral families fail strong summability
outside $U_q$. We did not prove that every possible resummation,
noncanonical assignment, or change of analytic category is impossible.
Any proposed extension must specify its summation rule and which
identities it preserves; it cannot call the original families strongly
Hahn-summable when their leading supports are decreasing.

\paragraph{Independent verification and priority.}
The most important independent proof checks are the joint-family
interchange in \cref{lem:agreement} and the use of regular projective
charts in the group law. The source search was not an exhaustive search
of valuation theory, formal models, or the model theory of analytic
structures. Finding an earlier theorem with these conclusions would
change the priority assessment, not the support calculation. The
article should be treated as a research draft with explicit proofs,
not as an independently certified publication.

\section{Conclusion}

At arbitrary valuation rank, the correct multiplicative source of the
Hahn Tate map is neither automatically all of $K^\times$ nor merely
one rank-one coefficient subfield. It is
\[
 U_q=\{u:v(u)\in H_{v(q)}\}.
\]
The restriction is forced by the bilateral support geometry, yet it
loses no elliptic-curve points. Classical rank-one uniformization reaches
the coarse residue curve; formal evaluation at coarse infinitesimals
recovers every remaining point. This proves
\[
 U_q/q^\Z\simeq E_q(K)
\]
without claiming convergence at scales where the series do not converge.

The same analysis yields a canonical valuation quotient, exact behavior
under scalar extension, support-preserving modular inversion, and an
explicit torsion ramification law. In particular, the construction gives
a coherent theory of multiplicatively uniformized elliptic curves over
$\No[i]$ while retaining its genuinely non-Archimedean scale structure.

\appendix
\section{Source and verification audit}\label{app:audit}

\subsection{What is attributed, and what is proved here}

The principal classical inputs are Tate's complete rank-one
uniformization theorem, the formal Tate curve and its discriminant,
Jacobi's product, algebraic closedness of divisible Hahn fields over
$\C$, and standard finite torsion and $j$-classification of elliptic
curves. They are cited rather than presented as discoveries.

The proof supplied here establishes the exact strong-summability domains,
the support and precision certificates, unrestricted formal evaluation
across a convex coarsening, its agreement with the literal Tate sums,
the resulting arbitrary-rank surjective quotient, and the indicated
consequences. The finite computations independently check selected
identities and signs. They neither replace the classical inputs nor
certify the all-rank proof.

\subsection{Repository provenance}

The reference revision identified during the audit was
\begin{center}\small\sha.\end{center}
The rank-one README was fetched explicitly at that revision. The root
README, document map, formalization ledger, and spectral README were
read from \texttt{main} during the same audit. No claim is made that all
manuscripts in the repository were read line by line. In particular,
the negative result of a keyword search is treated as a navigation aid,
not a formal proof that a topic occurs nowhere.

The delivered files do not alter the repository or claim to be part of
its checked Lean build. They form a standalone research package.

\subsection{Literature-search scope}

The audit checked Tate's original complete-field formulation and
subsequent treatments; the explicit formulas in the 2026
Eterovi\'c--Vermeulen paper; rank-one skeleton statements in
Baker--Payne--Rabinoff; and the existence and scope of higher-rank Hahn
analytification in Foster--Ranganathan. Searches combined Tate
uniformization, Hahn fields, higher-rank valuations, and theta series.
No matching exact-domain/all-points statement was identified in that
limited search. This is not a claim that such a result is absent from
all published or unpublished literature.

\subsection{Reproduction}

The source uses standard TeX Live packages and requires no external
figures, bibliography download, or font installation. From the package
root, run
\begin{verbatim}
python code/check_identities.py
latexmk -pdf -interaction=nonstopmode -halt-on-error article.tex
\end{verbatim}
The first command requires SymPy. The recorded verification output is
included so that the finite results remain inspectable without running
Python. Recompiling the article does not execute the verification code.

\begin{thebibliography}{99}\small

\bibitem{Tate}
J.~Tate,
\emph{A review of non-Archimedean elliptic functions},
in J.~Coates and S.-T.~Yau (eds.),
\emph{Elliptic Curves, Modular Forms, \& Fermat's Last Theorem},
International Press, 1995, pp.~162--184.
Author manuscript, including the 1959 proof:
\url{https://web.ma.utexas.edu/users/voloch/Preprints/nonarch-ams.pdf}.

\bibitem{Silverman}
J.~H.~Silverman,
\emph{Advanced Topics in the Arithmetic of Elliptic Curves},
Graduate Texts in Mathematics 151, Springer, 1994, Chapter~V.

\bibitem{SilvermanArithmetic}
J.~H.~Silverman,
\emph{The Arithmetic of Elliptic Curves},
2nd ed., Graduate Texts in Mathematics 106, Springer, 2009.

\bibitem{BGR}
S.~Bosch, U.~G\"untzer, and R.~Remmert,
\emph{Non-Archimedean Analysis}, Grundlehren der mathematischen
Wissenschaften 261, Springer, 1984.

\bibitem{Conrad}
B.~Conrad,
\emph{Math 248B: Algebraic theory of $q$-expansions}, course handout,
Stanford University.
\url{https://math.stanford.edu/~conrad/248BPage/handouts/qexp.pdf}.

\bibitem{EV}
S.~Eterovi\'c and F.~Vermeulen,
\emph{Hensel minimality, $p$-adic exponentiation and Tate uniformization},
arXiv:2602.16433, 2026, version 1, especially Sections~2.2 and~3.
\url{https://arxiv.org/abs/2602.16433}.

\bibitem{BPR}
M.~Baker, S.~Payne, and J.~Rabinoff,
\emph{On the structure of nonarchimedean analytic curves},
arXiv:1404.0279, 2014, especially Remark~4.24.
\url{https://arxiv.org/abs/1404.0279}.

\bibitem{FR}
T.~Foster and D.~Ranganathan,
\emph{Hahn analytification and connectivity of higher rank tropical
varieties}, arXiv:1504.07207, 2015.
\url{https://arxiv.org/abs/1504.07207}.

\bibitem{Neumann}
B.~H.~Neumann,
\emph{On ordered division rings},
Transactions of the American Mathematical Society \textbf{66} (1949),
202--252.

\bibitem{Higman}
G.~Higman,
\emph{Ordering by divisibility in abstract algebras},
Proceedings of the London Mathematical Society (3) \textbf{2} (1952),
326--336.

\bibitem{Poonen}
B.~Poonen,
\emph{Maximally complete fields},
L'Enseignement Math\'ematique (2) \textbf{39} (1993), 87--106.
\url{https://math.mit.edu/~poonen/papers/amsval.pdf}.

\bibitem{Gonshor}
H.~Gonshor,
\emph{An Introduction to the Theory of Surreal Numbers},
London Mathematical Society Lecture Note Series 110,
Cambridge University Press, 1986.

\bibitem{BournezGuilmant}
O.~Bournez and Q.~Guilmant,
\emph{Surreal fields stable under exponential and logarithmic
functions}, arXiv:2201.08199, 2022.
\url{https://arxiv.org/abs/2201.08199}.

\bibitem{DLMF}
F.~W.~J.~Olver et al. (eds.),
\emph{NIST Digital Library of Mathematical Functions},
Chapter~20, especially Section~20.5, infinite products for theta
functions. Accessed 22 September 2026.
\url{https://dlmf.nist.gov/20.5}.

\bibitem{Repo}
V.~Reshetnikov, \emph{Surreal}, GitHub repository,
root README, document map, and formalization ledger;
audit reference revision \texttt{a3124af}.
Accessed 22 September 2026.
\url{https://github.com/VladimirReshetnikov/Surreal}.

\bibitem{RepoRankOne}
\emph{Rank-One Non-Archimedean Analytic Geometry inside the Surcomplex
Numbers: Tate algebras, Berkovich disks, analytic annuli, and an
entire-function case study}, research package in \cite{Repo},
\texttt{docs/surcomplex/rank-one-berkovich/}.
README inspected at revision \texttt{a3124af}.
This is a repository research draft, not an independently refereed source.

\end{thebibliography}
\end{document}
